% GENERATED FILE. Edit canonical lessons, then run npm run generate:books.
% canonical-source-hash: fnv1a64:f99c68c789186123
% canonical-lessons: ML-C56-mat, ML-C56-basket, ML-C56-knife, ML-C56-vessel, ML-C56-box

\chapter{Things About the House}
\label{ch:ml-things-about-the-house}
\hlchaptermodality{56}

\begin{chapteropening}
\textbf{By the end of this chapter:} \emph{I can name five ordinary household things in Malayalam.}

Complete ML-C56-box and say all five words in order.
\end{chapteropening}

\section[pāya]{\ml{പായ} (pāya) — a mat}

\label{lesson:ML-C56-mat}



\begin{quote}
Before the new run of words: what did \emph{vaḻi} mean, and what did \emph{vayal} mean?
\end{quote}

\subsection*{You'll want to know: \ml{പായ}}
\textbf{\ml{പായ}} (\emph{pāya}) — \textquotedblleft{}a mat\textquotedblright{}.

Inherited, matching Tamil \emph{pāy}, and it names the woven mat that a Kerala household spreads to sleep on and unrolls for a guest to sit on.

The material is usually screwpine leaf or reed, split and plaited, and the making of them is a village trade with its own vocabulary. A \ml{പായ} rolls up and stands in a corner, which is why a house with few chairs can still seat twenty people.

The mat and the chair answer the same question by different means, and the older answer is the one still in daily use.

\begin{grammarlens}[title={What you've built}]
The first of five things a Kerala house keeps.
\end{grammarlens}

\subsection*{Your turn}
Say these aloud:

\begin{itemize}
\raggedright
  \item \emph{pāya}
  \item it once more, slowly
  \item \emph{vayal}, then \emph{pāya}, and say which of the two is woven
\end{itemize}

\begin{tcolorbox}[breakable,colback=teal!4,colframe=teal!35!black,title={Before you move on}]
What does \ml{പായ} mean? (\textquotedblleft{}a mat\textquotedblright{}.) Say it once more.
\end{tcolorbox}
\section[koṭṭa]{\ml{കൊട്ട} (koṭṭa) — a basket}

\label{lesson:ML-C56-basket}



\begin{quote}
Before the new one: what did \emph{pāya} mean?
\end{quote}

\subsection*{You'll want to know: \ml{കൊട്ട}}
\textbf{\ml{കൊട്ട}} (\emph{koṭṭa}) — \textquotedblleft{}a basket\textquotedblright{}.

Inherited, and plaited from bamboo or palm strip like the mat before it. A \ml{കൊട്ട} is the wide open basket that carries fish up from a boat, rice out of a field, or a load of pepper to market.

Kerala sizes them by trade rather than by measure, so a fish \ml{കൊട്ട} and a grain \ml{കൊട്ട} are different objects under one word, and the market knows which is meant.

The doubled \ml{ട} in the middle is held, the way the doubled sound in the seed-word was. Malayalam does a great deal of work with the length of a consonant.

\begin{grammarlens}[title={What you've built}]
Two, and both of them plaited.
\end{grammarlens}

\subsection*{Your turn}
Say these aloud:

\begin{itemize}
\raggedright
  \item \emph{koṭṭa}
  \item it once more, slowly
  \item \emph{pāya}, then \emph{koṭṭa}, and hold the doubled sound in the second
\end{itemize}

\begin{tcolorbox}[breakable,colback=teal!4,colframe=teal!35!black,title={Before you move on}]
What does \ml{കൊട്ട} mean? (\textquotedblleft{}a basket\textquotedblright{}.) Say it once more.
\end{tcolorbox}
\section[katti]{\ml{കത്തി} (katti) — a knife}

\label{lesson:ML-C56-knife}



\begin{quote}
Before the new one: what did \emph{koṭṭa} mean?
\end{quote}

\subsection*{You'll want to know: \ml{കത്തി}}
\textbf{\ml{കത്തി}} (\emph{katti}) — \textquotedblleft{}a knife\textquotedblright{}.

Inherited, matching Tamil \emph{katti}, and the doubled \emph{tt} is held long here too — three words into this run and the pattern is doing its own teaching.

Kerala's characteristic blade is not held in the hand at all. It is the \ml{വാക്കത്തി} and the sickle-shaped cutter clamped to a board, worked by pulling the vegetable across a fixed edge while sitting on the floor. The word covers that as readily as it covers a table knife.

A word for a tool tends to outlast the tool's shape, and this one has outlasted several.

\begin{grammarlens}[title={What you've built}]
Three: something to sit on, something to carry with, something to cut with.
\end{grammarlens}

\subsection*{Your turn}
Say these aloud:

\begin{itemize}
\raggedright
  \item \emph{katti}
  \item it once more, slowly
  \item \emph{koṭṭa}, then \emph{katti}, and hear the same held sound in both
\end{itemize}

\begin{tcolorbox}[breakable,colback=teal!4,colframe=teal!35!black,title={Before you move on}]
What does \ml{കത്തി} mean? (\textquotedblleft{}a knife\textquotedblright{}.) Say it once more.
\end{tcolorbox}
\section[pātraṁ]{\ml{പാത്രം} (pātraṁ) — a vessel, a container}

\label{lesson:ML-C56-vessel}



\begin{quote}
Before the new one: what did \emph{katti} mean?
\end{quote}

\subsection*{You'll want to know: \ml{പാത്രം}}
\textbf{\ml{പാത്രം}} (\emph{pātraṁ}) — \textquotedblleft{}a vessel, a container\textquotedblright{}.

Sanskrit \emph{pātra}, \textquotedblleft{}a receptacle\textquotedblright{} — the first borrowed word in this run, and it arrived carrying a second sense that Malayalam kept.

A \ml{പാത്രം} is a cooking vessel, a pot, a dish. It is also a character in a story or a part in a play: the person is a vessel that a role is poured into. Malayalam film writing uses the word constantly in that second sense, and a reader meets both without difficulty.

One shape holding the pot and the person is the kind of figure a language borrows whole rather than builds, and Malayalam took it as it stood.

\begin{grammarlens}[title={What you've built}]
Four, and this one holds a story as well as a stew.
\end{grammarlens}

\subsection*{Your turn}
Say these aloud:

\begin{itemize}
\raggedright
  \item \emph{pātraṁ}
  \item it once more, slowly
  \item \emph{pātraṁ}, and then say which of its two senses a cook means
\end{itemize}

\begin{tcolorbox}[breakable,colback=teal!4,colframe=teal!35!black,title={Before you move on}]
What does \ml{പാത്രം} mean? (\textquotedblleft{}a vessel, a container\textquotedblright{}.) Say it once more.
\end{tcolorbox}
\section[peṭṭi]{\ml{പെട്ടി} (peṭṭi) — a box}

\label{lesson:ML-C56-box}



\begin{quote}
Before the new one: what did \emph{pātraṁ} mean?
\end{quote}

\subsection*{You'll want to know: \ml{പെട്ടി}}
\textbf{\ml{പെട്ടി}} (\emph{peṭṭi}) — \textquotedblleft{}a box\textquotedblright{}.

Matching Tamil \emph{peṭṭi}, and it covers the whole family of closed containers: a box, a trunk, a case, a crate.

It is one of those words that quietly absorbs new objects. Malayalam has built the matchbox, the ballot box and the letter box on it since, each one a compound with \ml{പെട്ടി} sitting on the end and an older word in front.

A word that keeps accepting new referents without changing shape is a sign of a healthy one, and \ml{പെട്ടി} has been accepting them for a long time.

\begin{grammarlens}[title={What you've built}]
Five, and the house is furnished: a mat, a basket, a knife, a vessel, a box.
\end{grammarlens}

\subsection*{Your turn}
Say these aloud:

\begin{itemize}
\raggedright
  \item \emph{peṭṭi}
  \item it once more, slowly
  \item all five in order — \emph{pāya}, \emph{koṭṭa}, \emph{katti}, \emph{pātraṁ}, \emph{peṭṭi}
\end{itemize}

\begin{tcolorbox}[breakable,colback=teal!4,colframe=teal!35!black,title={Before you move on}]
What does \ml{പെട്ടി} mean? (\textquotedblleft{}a box\textquotedblright{}.) Say it once more.
\end{tcolorbox}
